\documentclass[11pt,a4paper]{article}
\usepackage{fontspec,amsmath,amssymb,amsthm,microtype}
\setmainfont{lmroman10-regular.otf}[BoldFont=lmroman10-bold.otf,ItalicFont=lmroman10-italic.otf,BoldItalicFont=lmroman10-bolditalic.otf]
\usepackage{xcolor}
\usepackage[most]{tcolorbox}
\usepackage{mathrsfs}
\providecommand{\lt}{<}
\providecommand{\gt}{>}
\setmonofont{DejaVuSansMono.ttf}[BoldFont=DejaVuSansMono-Bold.ttf,ItalicFont=DejaVuSansMono-Oblique.ttf,BoldItalicFont=DejaVuSansMono-BoldOblique.ttf,Scale=0.86]
\newfontfamily\leanfallback{STIXTwoMath-Regular.otf}[Scale=0.86]
\definecolor{leanrule}{gray}{0.62}
\definecolor{leanback}{gray}{0.962}
\definecolor{leantitle}{gray}{0.22}
\definecolor{leanmuted}{gray}{0.38}
\newtcolorbox{leanbox}[1]{enhanced,breakable,colback=leanback,colframe=leanrule,
  boxrule=0.4pt,arc=0.6mm,left=2.4mm,right=2.4mm,top=1.8mm,bottom=1.8mm,
  before skip=0.8em,after skip=1em,fonttitle=\small\bfseries,coltitle=leantitle,
  colbacktitle=leanback,titlerule=0.3pt,title={#1}}
\newlength{\leancharwidth}
\newenvironment{leancode}{\par\smallskip\begingroup\ttfamily\footnotesize\raggedright
  \settowidth{\leancharwidth}{M}\setlength{\parskip}{0pt}\setlength{\parindent}{0pt}
  \frenchspacing}{\par\endgroup\smallskip}
\newcommand{\leanline}[2]{\hangindent=\dimexpr#1\leancharwidth+4\leancharwidth\relax
  \hangafter=1\noindent\hspace*{#1\leancharwidth}#2\par}
\newcommand{\leansource}[1]{\par{\footnotesize\color{leanmuted}#1}\par}
\newcommand{\leanlabel}[1]{\par\smallskip\noindent{\small\bfseries #1}\par}
\usepackage[margin=25mm]{geometry}
\usepackage[unicode,hidelinks]{hyperref}
\hypersetup{pdfauthor={LangevinLean Project}}
\setlength{\emergencystretch}{3em}
\setlength{\parskip}{0.4em}
\setlength{\parindent}{0pt}
\allowdisplaybreaks[2]
\raggedbottom
\title{\texorpdfstring{Differences and uniqueness of common-input integral solutions}{Differences and uniqueness of common-input integral solutions}}
\author{LangevinLean Project}
\date{}
\begin{document}\maketitle
Draft mathematical exposition; these notes may contain mistakes.


Individual kinetic velocity paths contain continuous noise and need not be differentiable. Their difference under the same additive forcing is differentiable from the right, because the forcing cancels in the integral equations themselves. This section derives that assertion and proves uniqueness by Gronwall, including the initial time.



Cancellation is performed at the level of integral equations, where Brownian differentiability is unnecessary. The resulting deterministic difference equation supports both uniqueness and the transformed contraction theorem.


\section{\texorpdfstring{Subtract the two integral equations}{Subtract the two integral equations}}


Suppose \(Y_t-Y_0=-\int_0^t g(Y_s)ds+B_t\) and \(Z_t-Z_0=-\int_0^t g(Z_s)ds+B_t\), with continuous paths and a Lipschitz drift. Subtraction gives

\[Y_t-Z_t=Y_0-Z_0-\int_0^t[g(Y_s)-g(Z_s)]ds.\]

The two drift integrals exist on every compact interval because their integrands are continuous. First, the argument establishes interval integrability, distributes the difference through the integral, rewrites both supplied equations as expressions for the paths, and cancels the common forcing with the additive-group calculation. No derivative of the noise is taken.



The initial phase difference remains after subtraction, but the common forcing cancels identically at every time. The difference is an integral of the drift difference. This cancellation is valid for every continuous forcing path and does not use a distributional property of the noise.

\begin{leanbox}{Lean: \texttt{integral\_\allowbreak{}difference\_\allowbreak{}identity}}
\leansource{\texttt{integral\_\allowbreak{}difference\_\allowbreak{}identity} \quad KineticLangevin/\allowbreak{}IntegralCoupling.lean, lines 21–37}
\begin{leancode}
\leanline{0}{theorem\ integral\_difference\_identity}
\leanline{4}{(g\ :\ E\ →\ E)\ (L\ :\ ℝ≥0)\ (hg\ :\ LipschitzWith\ L\ g)}
\leanline{4}{(Y\ Z\ B\ :\ ℝ\ →\ E)\ (hY\ :\ Continuous\ Y)\ (hZ\ :\ Continuous\ Z)\ (T\ :\ ℝ)}
\leanline{4}{(heY\ :\ ∀\ t\ ∈\ Icc\ 0\ T,\ Y\ t\ -\ Y\ 0\ =\ -(∫\ s\ in\ 0..t,\ g\ (Y\ s))\ +\ B\ t)}
\leanline{4}{(heZ\ :\ ∀\ t\ ∈\ Icc\ 0\ T,\ Z\ t\ -\ Z\ 0\ =\ -(∫\ s\ in\ 0..t,\ g\ (Z\ s))\ +\ B\ t)}
\leanline{4}{(t\ :\ ℝ)\ (ht\ :\ t\ ∈\ Icc\ 0\ T)\ :}
\leanline{4}{Y\ t\ -\ Z\ t\ =\ Y\ 0\ -\ Z\ 0\ -\ ∫\ s\ in\ 0..t,\ g\ (Y\ s)\ -\ g\ (Z\ s)}
\end{leancode}
\leanlabel{How the Lean proof follows the argument}
\texttt{integral\_\allowbreak{}difference\_\allowbreak{}identity} gives the formal counterpart of this step. The main intermediate facts are \texttt{hi\allowbreak{}Y}, \texttt{hi\allowbreak{}Z}, \texttt{hy}, \texttt{hz} in that order.
\end{leanbox}

\section{\texorpdfstring{Differentiate after cancellation, including the initial endpoint}{Differentiate after cancellation, including the initial endpoint}}


The integrand \(g(Y_s)-g(Z_s)\) is continuous. The fundamental theorem of calculus therefore differentiates the integral expression, producing the right derivative \(-(g(Y_t)-g(Z_t))\). The identity is known only for \(t\in[0,T]\); consequently the mathematically appropriate conclusion at zero is a derivative within \([t,\infty)\). The proof builds an eventual neighborhood lying inside the valid time interval and transfers the derivative from the integral formula using eventual equality. This avoids imposing an arbitrary differentiable extension of the stochastic path to negative time.



The integrand is continuous because the force and the two solution paths are continuous. The fundamental theorem of calculus therefore gives an ordinary derivative in the interior and a right derivative at zero. These are exactly the derivatives used by the transformed-energy calculation.

\begin{leanbox}{Lean: \texttt{integral\_\allowbreak{}difference\_\allowbreak{}hasDerivWithinAt}}
\leansource{\texttt{integral\_\allowbreak{}difference\_\allowbreak{}hasDerivWithinAt} \quad KineticLangevin/\allowbreak{}IntegralCoupling.lean, lines 41–61}
\begin{leancode}
\leanline{0}{theorem\ integral\_difference\_hasDerivWithinAt}
\leanline{4}{(g\ :\ E\ →\ E)\ (L\ :\ ℝ≥0)\ (hg\ :\ LipschitzWith\ L\ g)}
\leanline{4}{(Y\ Z\ B\ :\ ℝ\ →\ E)\ (hY\ :\ Continuous\ Y)\ (hZ\ :\ Continuous\ Z)\ (T\ :\ ℝ)}
\leanline{4}{(heY\ :\ ∀\ t\ ∈\ Icc\ 0\ T,\ Y\ t\ -\ Y\ 0\ =\ -(∫\ s\ in\ 0..t,\ g\ (Y\ s))\ +\ B\ t)}
\leanline{4}{(heZ\ :\ ∀\ t\ ∈\ Icc\ 0\ T,\ Z\ t\ -\ Z\ 0\ =\ -(∫\ s\ in\ 0..t,\ g\ (Z\ s))\ +\ B\ t)}
\leanline{4}{(t\ :\ ℝ)\ (ht\ :\ t\ ∈\ Ico\ 0\ T)\ :}
\leanline{4}{HasDerivWithinAt\ (fun\ r\ ↦\ Y\ r\ -\ Z\ r)\ (-(g\ (Y\ t)\ -\ g\ (Z\ t)))\ (Ici\ t)\ t}
\end{leancode}
\leanlabel{How the Lean proof follows the argument}
\texttt{integral\_\allowbreak{}difference\_\allowbreak{}has\allowbreak{}Deriv\allowbreak{}Within\allowbreak{}At} gives the formal counterpart of this step. The main intermediate facts are \texttt{hforce}, \texttt{hd}, \texttt{hmem}, \texttt{he} in that order. Named proof applications include \texttt{\ensuremath{\langle}ht.\allowbreak{}1.\allowbreak{}trans}, \texttt{integral\_\allowbreak{}difference\_\allowbreak{}identity}; the adjacent proof shows their hypotheses and order of use.
\end{leanbox}

\section{\texorpdfstring{Uniqueness from the global Lipschitz bound}{Uniqueness from the global Lipschitz bound}}


For \(D=Y-Z\), the derived derivative obeys \(\|D\prime_t\|\le L\|D_t\|\). Equal initial conditions give \(\|D_0\|=0\). Gronwall then yields \(\|D_t\|\le0\) throughout the finite horizon. The proof passes the right-derivative theorem to its norm version of Gronwall, simplifies the zero initial bound, and uses definiteness of the norm to recover \(Y_t=Z_t\). This result supplies uniqueness for any continuous additive forcing, including the forcing used in the kinetic Picard construction.



Let \(D(t)\) be the norm of the difference of two solutions with the same initialization and input. Global Lipschitz continuity gives \(D(t)\le L\int_0^tD(s)ds\). Grönwall yields \(D(t)=0\), so every selected solution agrees with every other one on the horizon.

\begin{leanbox}{Lean: \texttt{integral\_\allowbreak{}solution\_\allowbreak{}unique}}
\leansource{\texttt{integral\_\allowbreak{}solution\_\allowbreak{}unique} \quad KineticLangevin/\allowbreak{}IntegralCoupling.lean, lines 65–80}
\begin{leancode}
\leanline{0}{theorem\ integral\_solution\_unique}
\leanline{4}{(g\ :\ E\ →\ E)\ (L\ :\ ℝ≥0)\ (hg\ :\ LipschitzWith\ L\ g)}
\leanline{4}{(Y\ Z\ B\ :\ ℝ\ →\ E)\ (hY\ :\ Continuous\ Y)\ (hZ\ :\ Continuous\ Z)\ (T\ :\ ℝ)}
\leanline{4}{(heY\ :\ ∀\ t\ ∈\ Icc\ 0\ T,\ Y\ t\ -\ Y\ 0\ =\ -(∫\ s\ in\ 0..t,\ g\ (Y\ s))\ +\ B\ t)}
\leanline{4}{(heZ\ :\ ∀\ t\ ∈\ Icc\ 0\ T,\ Z\ t\ -\ Z\ 0\ =\ -(∫\ s\ in\ 0..t,\ g\ (Z\ s))\ +\ B\ t)}
\leanline{4}{(hinitial\ :\ Y\ 0\ =\ Z\ 0)\ (t\ :\ ℝ)\ (ht\ :\ t\ ∈\ Icc\ 0\ T)\ :\ Y\ t\ =\ Z\ t}
\end{leancode}
\leanlabel{How the Lean proof follows the argument}
\texttt{integral\_\allowbreak{}solution\_\allowbreak{}unique} gives the formal counterpart of this step. The main intermediate facts are \texttt{hd}, \texttt{hgrowth}, \texttt{hb} in that order. Named proof applications include \texttt{hg.\allowbreak{}norm\_\allowbreak{}sub\_\allowbreak{}le}, \texttt{sub\_\allowbreak{}eq\_\allowbreak{}zero.\allowbreak{}mp}; the adjacent proof shows their hypotheses and order of use.
\end{leanbox}

\end{document}
